\documentclass[twoside,letterpaper,11pt]{article}
\usepackage{lshort2e,mylayout}

\begin{document}

\title{Essential \LaTeXe\ i18n}
\author{Tobias Oetiker}
\date{2002}

\maketitle

\index{international} When you write documents in \wi{language}s
other than English, there are three areas where \LaTeX{} has to be
configured appropriately:

\begin{enumerate}
\item All automatically generated text strings\footnote{Table of
    Contents, List of Figures, \ldots} have to be adapted to the new
  language.  For many languages, these changes can be accomplished by
  using the \pai{babel} package by Johannes Braams.
\item \LaTeX{} needs to know the hyphenation rules for the new
  language. Getting hyphenation rules into \LaTeX{} is a bit more
  tricky. It means rebuilding the format file with different
  hyphenation patterns enabled. Your \guide{} should give more
  information on this.
\item Language specific typographic rules. In French for example, there is a
  mandatory space before each colon character (:).
\end{enumerate}

\section{The Babel Package}

If your system is already configured appropriately, you can activate
the \pai{babel} package by adding the command
\begin{lscommand}
\ci{usepackage}\verb|[|\emph{language}\verb|]{babel}| 
\end{lscommand}
\noindent after the \verb|\documentclass| command. A list of the
\emph{language}s built into your \LaTeX{} system will be displayed
every time the compiler is started. Babel will
automatically activate the appropriate hyphenation rules for the
language you choose. If your \LaTeX{} format does not support
hyphenation in the language of your choice, babel will still work but
will disable hyphenation, which has quite a negative effect on the
appearance of the typeset document.

\textsf{Babel} also specifies new commands for some languages, which
simplify the input of special characters. The \wi{German} language, for
example, contains a lot of umlauts (\"a\"o\"u).  With \textsf{babel},
you can enter an \"o by typing \verb|"o| instead of~\verb|\"o|.

If you call babel with multiple languages
\begin{lscommand}
\ci{usepackage}\verb|[|\emph{languageA}\verb|,|\emph{languageB}\verb|]{babel}| 
\end{lscommand}
\noindent you have to use the command
\begin{lscommand}
\ci{selectlanguage}\verb|{|\emph{languageA}\verb|}|
\end{lscommand}
\noindent to set the current language.

\section{Input Encodings}

Most modern computer systems allow you to input some special characters
directly from the keyboard. \LaTeX{} can handle such characters through the
\pai{inputenc} package:
\begin{lscommand}
\ci{usepackage}\verb|[|\emph{encoding}\verb|]{inputenc}| 
\end{lscommand}

When using this package, you should consider
that other people might not be able to display your input files on
their computer, because they use a different encoding. For example,
the German umlaut \"a on OS/2 is encoded as 132, but on Unix
systems using ISO-LATIN~1 it is encoded as 228; therefore you should
use this feature with care. The following encodings may come in handy,
depending on the type of system you are working on: 

\begin{center}
\begin{tabular}{l | r}
Operating system & encoding\\
\hline
Mac     &  \texttt{applemac} \\
Unix    &  \texttt{latin1} \\ 
Windows &  \texttt{ansinew} \\
DOS, OS/2    &  \texttt{cp850}
\end{tabular}
\end{center}

\section{Font Encodings}

Font encoding is a different matter. It defines at which position inside
a \TeX-font each letter is stored. The original Computer Modern
\TeX{} font only contains the 128 characters of the old 7-bit ASCII
character set. When accented characters are required, \TeX{} creates
them by combining a normal character with an accent. While the
resulting output looks perfect, this approach stops the automatic
hyphenation from working inside words containing accented characters.

Fortunately, most modern \TeX{} distributions contain a copy of the EC
fonts. These fonts look like the Computer Modern fonts, but contain
special characters for most of the accented characters used in
European languages. By using these fonts you can improve hyphenation
in non-English documents. The EC fonts are activated by including the
\pai{fontenc} package in the preamble of your document. \label{fontenc}

\begin{lscommand}
\ci{usepackage}\verb|[T1]{fontenc}| 
\end{lscommand}

\iffalse

As of this writing, no free, high quality \textsc{PostScript} versions of the EC fonts are
available.\footnote{Vladimir Volovich has created the cm-super font bundle
which covers the entire EC/TC, EC Concrete, EC Bright and LH fonts. It is avalable from \texttt{CTAN:/fonts/ps-type1/cm-super} and is
included with \TeX{}Live7 and Mik\TeX. Unfortunately these fonts are not
of the same typographic quality as the normal Type1 renderings of the CM fonts} This causes problems when generating PDF versions of your
documents---the resulting PDF file will look abysmal on screen. You might
want to have a look at the \pai{aeguill} package, aka \emph{Almost European
Computer Modern with Guillemets}, which uses the characters from the original Computer
Modern font, but rearranges them in the EC order and can thus use the
excellent type 1 format CM fonts available on most systems. There is also an \pai{ae} package, but it
does not provide proper rendering of French quotes, aka guillemets (see the section
on pdf\LaTeX{} on page \pageref{sec:pdftex}).

\fi

\subsection{Support for French}

\secby{Daniel Flipo}{daniel.flipo@univ-lille1.fr}
Some hints for those creating \wi{French} documents with \LaTeX{}:
you can load French language support with the command

\begin{lscommand}
\verb|\usepackage[frenchb]{babel}|
\end{lscommand}

Note that, for historical reasons, the name of \textsf{babel}'s option 
for French is either \emph{frenchb} or \emph{francais} but not \emph{french}.

This enables French hyphenation, if you have configured your
LaTeX system accordingly. It also changes all automatic text into
French: \verb+\chapter+ prints Chapitre, \verb+\today+ prints the current
date in French and so on. A set of new commands also
becomes available, which allows you to write French input files more
easily. Check out table \ref{cmd-french} for inspiration. 

\begin{table}[!htbp]
\caption{Special commands for French.} \label{cmd-french}
\begin{lined}{9cm}
\selectlanguage{french}
\begin{tabular}{ll}
\verb+\og guillemets \fg{}+         \quad &\og guillemets \fg \\[1ex]
\verb+M\up{me}, D\up{r}+            \quad &M\up{me}, D\up{r}  \\[1ex]
\verb+1\ier{}, 1\iere{}, 1\ieres{}+ \quad &1\ier{}, 1\iere{}, 1\ieres{}\\[1ex]
\verb+2\ieme{} 4\iemes{}+           \quad &2\ieme{} 4\iemes{}\\[1ex]
\verb+\No 1, \no 2+                 \quad &\No 1, \no 2   \\[1ex]
\verb+20~\degres C, 45\degres+      \quad &20~\degres C, 45\degres \\[1ex]
\verb+\bsc{M. Durand}+              \quad &\bsc{M.~Durand} \\[1ex]
\verb+\nombre{1234,56789}+          \quad &\nombre{1234,56789}
\end{tabular}
\selectlanguage{english}
\bigskip
\end{lined}
\end{table}

You will also notice that the layout of lists changes when switching to the
French language. For more information on what the \texttt{frenchb}
option of \textsf{babel} does and how you can customize its behaviour, run
\LaTeX{} on file \texttt{frenchb.dtx} and read the produced file
\texttt{frenchb.dvi}.

\subsection{Support for German}

Some hints for those creating \wi{German}\index{Deutsch}
documents with \LaTeX{}: you can load German language support with the
command

\begin{lscommand}
\verb|\usepackage[german]{babel}|
\end{lscommand}

This enables German hyphenation, if you have configured your
\LaTeX system accordingly. It also changes all automatic text into
German. Eg. ``Chapter'' becomes ``Kapitel.'' A set of new commands also
becomes available, which allows you to write German input files more quickly
even when you don't use the inputenc package. Check out table
\ref{german} for inspiration. With inputenc, all this become mute, but your 
text also is locked in a particular encoding world.

\begin{table}[!htbp]
\caption{German Special Characters.} \label{german}
\begin{lined}{8cm}
\selectlanguage{german}
\begin{tabular}{*2{ll}}
\verb|"a| & "a \hspace*{1ex} & \verb|"s| & "s \\[1ex]
\verb|"`| & "` & \verb|"'| & "' \\[1ex]
\verb|"<| or \ci{flqq} & "<  & \verb|">| or \ci{frqq} & "> \\[1ex]
\ci{flq} & \flq & \ci{frq} & \frq \\[1ex]
\ci{dq} & " \\
\end{tabular}
\selectlanguage{english}
\bigskip
\end{lined}
\end{table}

In German books you often find French quotation marks (\flqq guillemets\frqq).
German typesetters, however, use them differently. A quote in a German book
would look like \frqq this\flqq. In the German speaking part of Switzerland,
typesetters use \flqq guillemets\frqq~the same way the French do.

A major problem arises from the use of commands
like \verb+\flq+: If you use the OT1 font (which is the default font) the  
guillemets will look like the math symbol ``$\ll$'', which turns a typesetter's stomach.
T1 encoded fonts, on the other hand, do contain the required symbols. So if you are using this type
of quote, make sure you use the T1 encoding. (\verb|\usepackage[T1]{fontenc}|)

\iffalse

\subsection{Support for Portuguese}

\secby{Demerson Andre Polli}{polli@linux.ime.usp.br}
To enable hyphenation and change all automatic text to \wi{Portuguese},
\index{Portugu\^es} use the command:
\begin{lscommand}
\verb|\usepackage[portuguese]{babel}|
\end{lscommand}
Or if you are in Brazil, substitute the language for \wi{brazilian}.

As there are a lot of accents in Portuguese you might want to use
\begin{lscommand}
\verb|\usepackage[latin1]{inputenc}|
\end{lscommand}
to be able to input them correctly as well as 
\begin{lscommand}
\verb|\usepackage[T1]{fontenc}|
\end{lscommand}
to get the hyphenation right.

See table~\ref{portuguese} for the preamble you need to write in the
Portuguese language. Note that we are using the latin1 input encoding here,
so this will not work on a Mac or on DOS. Just use
the appropriate encoding for your system.

\begin{table}[btp]
\caption{Preamble forPortuguese documents.} \label{portuguese}
\begin{lined}{5cm}
\begin{verbatim}
\usepackage[portuges]{babel}
\usepackage[latin1]{inputenc}
\usepackage[T1]{fontenc}
\end{verbatim}
\bigskip
\end{lined}
\end{table}

\subsection[Support for Korean]{Support for Korean\footnotemark}\label{support_korean}%
\footnotetext{%
Considering a number of issues  Korean \LaTeX{} users
have to cope with.
This section was written by Karnes KIM on behalf of the
Korean lshort translation team. It  was translated into English
by SHIN Jungshik and shortened by Tobi Oetiker}

To use \LaTeX{} for typesetting  \wi{Korean}, 
we need to solve three problems: 

\begin{enumerate}
\item 
We must be able to 
edit \wi{Korean input files}.
Korean input files must be in plain text format, but because Korean
uses its own character set outside the
repertoire of US-ASCII, they will look rather strange with a normal ASCII editor.  The two most widely used encodings for
Korean text files are  EUC-KR and its upward compatible
extension used in Korean MS-Windows, CP949/Windows-949/UHC.
In these encodings each US-ASCII character represents its normal ASCII
character similar to other ASCII compatible encodings such as
ISO-8859-\textit{x}, EUC-JP,Shift\_JIS, and Big5. On the other hand, Hangul
syllables, Hanjas (Chinese characters as used in Korea), Hangul Jamos,
Hirakanas, Katakanas, Greek and Cyrillic characters and other
symbols and letters drawn from KS~X~1001 are represented by two
consecutive octets. The first has its MSB set.
Until mid the 1990's, it took a considerable amount of time and effort to
set up a Korean-capable environment under a non-localized (non-Korean)
operating system. 
You can skim through the now much-outdated \url{http://jshin.net/faq} to get 
a glimpse of what it was like to use Korean under non-Korean OS in mid-1990's.
These days all three major operating systems (Mac OS, Unix, Windows) come equipped
with pretty decent multilingual support and internationalization features
so that editing Korean text file is not so much of a problem anymore, even
on non-Korean operating systems.

\item \TeX{} and \LaTeX{} were originally written for
scripts with no more than 256 characters in their alphabet.
To make them work for languages with considerably 
more characters such as
Korean%,
 \footnote{Korean Hangul is an alphabetic script with 14 basic consonants
 and 10 basic vowels (Jamos). Unlike Latin or Cyrillic scripts, the
 individual characters have to be arranged in rectangular
 clusters about the same size as Chinese characters. Each cluster
 represents a syllable. An unlimited number of syllables can be
 formed out of this finite set of vowels and consonants. Modern Korean
 orthographic standards (both in South Korea and  North Korea), however,
 put some restriction on the formation of these clusters.
 Therefore only a finite the number of  orthographically correct syllables exist.
 The Korean Character encoding defines individual code points for each of these syllables (KS~X~1001:1998 and KS~X~1002:1992). So Hangul, albeit alphabetic, is
 treated like the Chinese and Japanese writing systems with tens of thousands of
 ideographic/logographic characters.  ISO~10646/Unicode offers both ways of
 representing Hangul used for \emph{modern} Korean by encoding Conjoining
 Hangul Jamos (alphabets: \url{http://www.unicode.org/charts/PDF/U1100.pdf})
 in addition to encoding all the orthographically allowed Hangul syllables in
 \emph{modern} Korean (\url{http://www.unicode.org/charts/PDF/UAC00.pdf}).
 One of the most daunting challenges in Korean typesetting with
 \LaTeX{} and related typesetting system is supporting Middle Korean---and possibly future Korean---syllables that can be only represented
 by conjoining Jamos in Unicode. It is hoped that future \TeX{} engines like $\Omega$ and
 $\Lambda$ will eventually provide solutions to this
 so that some Korean linguists and historians
 will defect from MS Word that already has  a pretty good support 
 for Middle Korean.}
or Chinese, a subfont mechanism was developed.
It divides a single CJK fonts with  thousands or tens of thousands of
glyphs into a set of subfonts with 256 glyphs each. 
For Korean, there are three widely used packages;  \wi{H\LaTeX}
by UN~Koaunghi, \wi{h\LaTeX{}p} by CHA~Jaechoon and the \wi{CJK package}
by Werner~Lemberg.\footnote{% 
They can be obtained at \texttt{CTAN:/tex-archive/language/korean/HLaTeX/}\\
   \texttt{CTAN:/tex-archive/language/korean/CJK/} and
   \texttt{http://knot.kaist.ac.kr/htex/}}
H\LaTeX{} and h\LaTeX{}p are specific to Korean and provide
Korean localization on top of the font support.
They both can process Korean input text files encoded in EUC-KR. H\LaTeX{} can
even process input files encoded in CP949/Windows-949/UHC and UTF-8
when used along with $\Lambda$, $\Omega$.

The CJK package is not specific to Korean. It can
process input files in UTF-8 as well as in various CJK encodings
including EUC-KR and CP949/Windows-949/UHC, it can be used to typeset documents with
multilingual content (especially Chinese, Japanese and Korean).
The CJK package has no Korean localization such as the one offered by H\LaTeX{} and it
does not come with as many special Korean fonts as H\LaTeX.

\item The ultimate purpose of using typesetting programs like \TeX{}
and \LaTeX{} is to get documents typeset in an `aesthetically' satisfying way.
Arguably the most important element in typesetting is  a set of
well-designed fonts. The H\LaTeX{} distribution
includes \index{Korean font!UHC font}UHC \textsc{PostScript} fonts 
of 10
different families and
Munhwabu\footnote{Korean Ministry of Culture}
fonts (TrueType) of 5 different families.
The CJK package works with a set of fonts used by earlier versions
of H\LaTeX{} and it can use Bitstream's cyberbit TrueType
font.
\end{enumerate}

To use the  H\LaTeX{} package for typesetting your Korean text, put the following
declaration into the preamble of your document:
\begin{lscommand}
\verb+\usepackage{hangul}+
\end{lscommand}

This command turns the Korean localization on. The headings
of chapters, sections, subsections, table of content and table of
figures are all translated into Korean and the formatting of the document
is changed to follow Korean conventions. 
The packages also provides automatic ``particle selection.''
In Korean, there are pairs of post-fix particles 
grammatically equivalent but different in form. Which 
of any given pair is correct depends on 
whether the preceding syllable ends with a  vowel or a consonant.
(It is a bit more complex than this, but this should give you
a good picture.)
Native Korean speakers have no problem picking the right particle, but
it cannot be determined which particle to use for references and other automatic
text that will change while you edit the document.
It 
takes a painstaking effort to place appropriate particles manually
every time you add/remove references or simply shuffle  parts
of your document around.
H\LaTeX{} relieves its users from this boring and error-prone process.

In case you don't need Korean localization features
but just want 
to  typeset Korean text, you can put the following line in the 
preamble, instead.
\begin{lscommand}
\verb+\usepackage{hfont}+
\end{lscommand}

For more details on typesetting  Korean with H\LaTeX{}, refer to
the \emph{H\LaTeX{} Guide}.  Check out the web site of the Korean
\TeX{} User Group (KTUG) at  \url{http://www.ktug.or.kr/}.
There is also a Korean translation
of this manual available.

\fi

\end{document}